% ===========================================================================
% Preamble: packages, theorem environments, macros.
% Style mirrors the Zinc+ paper and the F_2 companion notes
% (sha-f2x-doc, zip-plus-f2x-pcs-doc).
% ===========================================================================

\usepackage[a4paper,margin=1in]{geometry}
\usepackage[utf8]{inputenc}
\usepackage[T1]{fontenc}
\usepackage{lmodern}
\usepackage{microtype}

\usepackage{amsmath,amssymb,amsthm,mathtools}
\usepackage{booktabs}
\usepackage{array}
\usepackage{tabularx}
\usepackage{multirow}
\usepackage{enumitem}
\usepackage{xcolor}
\usepackage[colorlinks=true,
            linkcolor={teal!70!black},
            citecolor={teal!70!black},
            urlcolor={teal!70!black}]{hyperref}
\usepackage{thmtools}
\usepackage{cleveref}

% --- Theorem environments --------------------------------------------------
% Shared `theorem` numbering via thmtools' `sibling=theorem`, with cleveref
% still naming each reference by its true kind ("Lemma 3.3", "Definition 5.1").
\theoremstyle{plain}
\declaretheorem[name=Theorem, numberwithin=section]{theorem}
\declaretheorem[name=Lemma, sibling=theorem]{lemma}
\declaretheorem[name=Proposition, sibling=theorem]{proposition}
\declaretheorem[name=Corollary, sibling=theorem]{corollary}
\declaretheorem[name=Claim, sibling=theorem]{claim}
\declaretheorem[name=Observation, sibling=theorem]{observation}

\theoremstyle{definition}
\declaretheorem[name=Definition, sibling=theorem]{definition}
\declaretheorem[name=Construction, sibling=theorem]{construction}
\declaretheorem[name=Protocol, sibling=theorem]{protocol}
\declaretheorem[name=Example, sibling=theorem]{example}

\theoremstyle{remark}
\declaretheorem[name=Remark, sibling=theorem]{remark}

% --- Operators -------------------------------------------------------------
\DeclareMathOperator{\ROTR}{ROTR}
\DeclareMathOperator{\ROTL}{ROTL}
\DeclareMathOperator{\SHR}{SHR}
\DeclareMathOperator{\SHL}{SHL}
\DeclareMathOperator{\XOR}{XOR}
\DeclareMathOperator{\AND}{AND}
\DeclareMathOperator{\OR}{OR}
\DeclareMathOperator{\NOT}{NOT}
\DeclareMathOperator{\Maj}{Maj}
\DeclareMathOperator{\Ch}{Ch}
\DeclareMathOperator{\MLE}{MLE}
\DeclareMathOperator{\eq}{eq}
\DeclareMathOperator{\enc}{Enc}
\DeclareMathOperator{\wt}{wt}
\DeclareMathOperator{\negl}{negl}
\DeclareMathOperator{\img}{img}
\DeclareMathOperator{\spanop}{span}
\DeclareMathOperator{\Merk}{Merkle}
\DeclareMathOperator{\polylog}{polylog}

% --- Number systems --------------------------------------------------------
\newcommand{\F}{\mathbb{F}}
\newcommand{\Ftwo}{\F_2}
\newcommand{\Q}{\mathbb{Q}}
\newcommand{\Z}{\mathbb{Z}}
\newcommand{\N}{\mathbb{N}}
\newcommand{\bP}{\mathbb{P}}

% System-name shorthands for the Binius64 comparison part.
\newcommand{\binius}{Binius64}
\newcommand{\zinc}{Zinc\texttt{+}}

% Binary extension field K = GF(2^kappa); deployed at kappa = 128.
\newcommand{\K}{\mathbb{K}}

% Cell ring R = F_2[X]/(X^w); deployed at word width w = 32.
\newcommand{\cellring}{R}
\newcommand{\wbit}{w}

% Bit-polynomial set: degree-<n with {0,1} coefficients.
\newcommand{\bitpoly}[1]{\{0,1\}^{<#1}[X]}

% Quotient rings used in the modular-addition arithmetisation.
\newcommand{\Fshift}{\Ftwo[X]/(X^{\wbit})}
\newcommand{\Frot}{\Ftwo[X]/(X^{\wbit}-1)}

% Degree-<d univariate over a coefficient ring S, in the cell variable X.
\newcommand{\polyat}[2]{#1[X]_{<#2}}
\newcommand{\Kpoly}{\polyat{\K}{\wbit}}        % K[X]_{<w}

% Projections (evaluation of the cell variable through an F_2-linear map).
\newcommand{\psia}{\psi_{\alpha}}
\newcommand{\psiz}{\psi_{z}}

% MLE tilde, Hadamard, bit-poly hat, ideal, column branding.
\newcommand{\mle}[1]{\widetilde{#1}}
\newcommand{\had}{\odot}
\newcommand{\bp}[1]{\widehat{#1}}
\newcommand{\ideal}[1]{\langle #1 \rangle}
\newcommand{\col}[1]{\ensuremath{\mathsf{#1}}}
\newcommand{\code}[1]{\ensuremath{\mathsf{#1}}}
\newcommand{\vv}[1]{\mathbf{#1}}
\newcommand{\mat}[1]{\mathsf{#1}}

% Hyperref-friendly section refs
\renewcommand{\sectionautorefname}{Section}
\renewcommand{\subsectionautorefname}{Section}
\renewcommand{\subsubsectionautorefname}{Section}

\usepackage{xcolor}

\newcommand{\albert}[1]{\textcolor{teal}{Albert: {#1}}}